\documentclass[a4paper]{article}
% Español (México) con XeLaTeX: fontspec para fuentes Unicode, polyglossia
% para la división silábica y los nombres del documento en la variante mexicana.
\usepackage{fontspec}
\usepackage{polyglossia}
\setdefaultlanguage[variant=mexican]{spanish}
% Los nombres chinos van como «pinyin (original)»: xeCJK da a los caracteres
% Han su propia fuente; sin ella XeLaTeX los omite con un aviso.
% FandolSong no cubre algunos caracteres de nombres (U+73FA); WenQuanYi Zen Hei sí.
\usepackage[AutoFallBack=true]{xeCJK}
\setCJKmainfont{FandolSong-Regular.otf}
\setCJKfallbackfamilyfont{\CJKrmdefault}{WenQuanYi Zen Hei}
% polyglossia no traduce los nombres de listings.
\AtBeginDocument{\providecommand{\lstlistingname}{}\renewcommand{\lstlistingname}{Listado}%
  \providecommand{\lstlistlistingname}{}\renewcommand{\lstlistlistingname}{Índice de listados}}
\usepackage{amsmath, amssymb}
\usepackage{graphicx}
\usepackage[margin=2.5cm]{geometry}
\usepackage[most]{tcolorbox}
\usepackage{etoolbox}
\usepackage{listings}
\usepackage{hyperref}
\usepackage{booktabs}
\usepackage{subcaption}
\usepackage{float}

\newtcolorbox{knowledgebox}[1]{enhanced, colback=blue!5!white, colframe=blue!75!black, colbacktitle=blue!75!black, coltitle=white, fonttitle=\bfseries, title={#1}, attach boxed title to top left={yshift=-2mm, xshift=2mm}, boxrule=1pt, sharp corners}
\newtcolorbox{importantbox}[1]{enhanced, colback=yellow!10!white, colframe=yellow!80!black, colbacktitle=yellow!80!black, coltitle=black, fonttitle=\bfseries, title={#1}, sharp corners}
\newtcolorbox{warningbox}[1]{enhanced, colback=red!5!white, colframe=red!75!black, colbacktitle=red!75!black, coltitle=white, fonttitle=\bfseries, title={#1}, sharp corners}

\lstset{language=Python, basicstyle=\ttfamily\small, keywordstyle=\color{blue}, stringstyle=\color{red!60!black}, commentstyle=\color{green!60!black}, breaklines=true, frame=single, numbers=left, numberstyle=\tiny\color{gray}, captionpos=b, extendedchars=true}

\newcommand{\notetitle}{{[}LLM Agents F24{]} Towards a Unified Framework of Neural and Symbolic Decision Making — Yuandong Tian}
\newcommand{\noteauthors}{Elaboradas a partir de materiales públicos del curso}
\newcommand{\notedate}{28 de octubre de 2024}
\newcommand{\videochannel}{Berkeley RDI}
\newcommand{\videopublishdate}{28 de octubre de 2024}
\newcommand{\videoduration}{1:04:58}
\newcommand{\videourl}{https://www.youtube.com/live/wm9-7VBpdEo}
\newcommand{\videocoverpath}{cover.jpg}
\newcommand{\repourl}{https://github.com/hqhq1025/ai-course-notes}

% Footer with repo info
\usepackage{fancyhdr}
\pagestyle{fancy}
\fancyhf{}
\fancyfoot[C]{\thepage}
\fancyfoot[R]{\footnotesize\color{gray}\href{\repourl}{AI Course Notes}}
\renewcommand{\headrulewidth}{0pt}
\renewcommand{\footrulewidth}{0pt}

\begin{document}
\begin{titlepage}
\centering
{\Large Notas de entrevista\par}
\vspace{1.2cm}
{\huge\bfseries \notetitle\par}
\vspace{0.8cm}
{\large \noteauthors\par}
\vspace{0.3cm}
{\large \notedate\par}
\vspace{1.2cm}

\ifdefempty{\videocoverpath}{
{\small Al generar las notas, indique la ruta local de la portada del video.\par}
}{
\includegraphics[width=0.82\textwidth,height=0.45\textheight,keepaspectratio]{\videocoverpath}\par
}

\vfill
\begin{tcolorbox}[width=0.9\textwidth, colback=black!2!white, colframe=black!60, sharp corners]
\textbf{Autor o canal del video}: \videochannel\par
\textbf{Fecha de publicación}: \videopublishdate\par
\textbf{Duración del video}: \videoduration\par
\ifdefempty{\videourl}{}{
\textbf{Enlace del video}: \href{\videourl}{\nolinkurl{\videourl}}\par
}
\textbf{Página del proyecto}: \href{\repourl}{\nolinkurl{\repourl}}\par
\end{tcolorbox}
\end{titlepage}

\tableofcontents
\newpage

\section{Contexto y desafíos del curso}

Esta clase de UC Berkeley CS294-196, otoño de 2024, la impartió Yuandong Tian, de Meta AI FAIR, con el título ``Towards a Unified Framework of Neural and Symbolic Decision Making''; se centra en cómo unificar las redes neuronales, la búsqueda, la planeación y el ciclo cerrado de datos de agentes en un mismo marco de decisión.
Abrió con una tendencia progresiva: a medida que aumenta la dificultad de la tarea, la curva de desempeño del mejor modelo sigue bajando de manera monótona, sin mostrar elasticidad.
\textit{``This is actually in some sense a very worrisome kind of curve that even the best model has not solved this planning problem.''}
Al observar repetidamente varios benchmarks (planeación inteligente, el juego Soang, la búsqueda de rutas en Maze), la misma tendencia aparece una y otra vez: la puntuación del modelo queda acotada por las propiedades simbólicas implícitas en el grafo de cómputo, y simplemente apilar parámetros no permite abrir el panorama en poco tiempo.
Esto marca el tono de las tres vías de respuesta de esta clase: depender solo de modelos más grandes y más costosos no basta.

\subsection{La curva en la capacidad de planeación}

En los primeros 10 minutos, el ponente mostró una curva de evaluación para tareas de planeación complejas: aunque la precisión del modelo puede mejorar, la puntuación tras la convergencia sigue siendo baja, y la cantidad de tokens o parámetros de entrenamiento necesarios crece de manera exponencial.
Este fenómeno se repite en diálogo, planeación automática y navegación robótica; la raíz del problema está en que un modelo generativo probabilístico no puede garantizar propiedades simbólicas (restricciones, factibilidad).
\begin{knowledgebox}{Síntomas del fracaso en la planeación}
Los síntomas más comunes de los LLM en la planeación incluyen: 1) violaciones de restricciones (por ejemplo, ignorar límites o restricciones de consumo de energía); 2) acumulación de errores en horizontes largos; 3) ausencia de una gestión precisa de las bifurcaciones del espacio de estados. Cualquier problema que dependa de la precisión, como la programación logística o la aprobación de cumplimiento normativo, se ve amplificado por esto.
\end{knowledgebox}

\subsection{Tres estrategias}

El ponente propuso tres vías de corrección: \textbf{Scaling} (aumentar tokens/parámetros), \textbf{Hybrid} (combinar LLM y solver) y, de manera más sistemática, \textbf{Compound AI}. No son opciones mutuamente excluyentes, sino una arquitectura organizada por capas según el costo y el alcance del razonamiento, y se pueden apilar una sobre otra.
\begin{enumerate}
    \item El extremo de \textbf{Scaling} confía en la generalización probabilística mediante un modelado del lenguaje más grande, más datos de entrenamiento y más cómputo, pero es poco sensible al conocimiento estructural de las variables de restricción.
    \item El extremo de \textbf{Hybrid} usa el LLM para aportar candidatos/acciones más heurísticas, mientras el solver se encarga de hacer cumplir las restricciones, reduciendo la falta de control del modelo generativo puro.
    \item El extremo de \textbf{Compound AI} encadena varios módulos (traza de búsqueda, diálogo de múltiples turnos, solver con restricciones) para formar un ciclo de retroalimentación auditable, lo que facilita la división del trabajo y la depuración.
\end{enumerate}
\begin{figure}[H]
\centering
\includegraphics[width=\textwidth]{frame-solution-categories.jpg}
\caption{El ponente enumera tres categorías de estrategias de solución, ubicadas en distintos niveles de costo de ingeniería y alcance de razonamiento. \protect\footnotemark}
\end{figure}
\footnotetext{Intervalo de tiempo en el video: 00:09:10--00:09:18.}
\begin{importantbox}{Desglose de las tres estrategias}
Scaling apuesta fuertemente por los recursos, Hybrid deja que el LLM aporte heurísticas mientras el solver se encarga de la búsqueda, y Compound AI encadena varios módulos para formar un ciclo de retroalimentación integrado, lo cual permite tanto controlar la calidad como facilitar el diagnóstico.
\end{importantbox}
\begin{warningbox}{No confíe ciegamente en Scaling}
Depender exclusivamente de modelos más grandes implica ciclos de entrenamiento más largos, sobreajuste más difícil de depurar y un consumo de energía en inferencia excesivo; si no se observa una mejora estructural, los parámetros adicionales terminarán únicamente por amplificar errores que ``parecen razonables''.
\end{warningbox}

\subsection{Resumen de la sección}

El mecanismo generativo de los LLM presenta una incertidumbre inherente en los problemas de planeación, y se necesita recurrir a estructuras simbólicas para recuperar el control; las tres estrategias ofrecen distintos puntos de entrada, tanto de ingeniería como metodológicos.

\section{Trayectorias de búsqueda y modelado potenciado por búsqueda}

Yuandong Tian continúa desarrollando esta línea central; el primer paso es tratar el propio proceso de búsqueda como una señal de supervisión, lo cual conserva la capacidad generativa del modelo neuronal y a la vez le da a los resultados una estructura interpretable.

\subsection{Trayectorias de búsqueda como weak supervision}

El expositor mostró la curva de entrenamiento del modelo potenciado por búsqueda: dentro del mismo orden de magnitud, el modelo de trayectorias de búsqueda alcanza un rendimiento superior al 80\%, mientras que el modelo puramente generativo se mantiene alrededor del 20\%. Esto se debe a que los tokens de salida del primero no solo contienen el plan final, sino también información como el elaborated trace y evaluaciones de cost, lo cual le brinda al Transformer más contexto.
\begin{figure}[H]
\centering
\includegraphics[width=\textwidth]{frame-search-curve.jpg}
\caption{Training curve del modelo potenciado por búsqueda comparado con el modelo puramente generativo, contrastando el baseline de 175M parámetros con el trace model de 15M parámetros. \protect\footnotemark}
\end{figure}
\footnotetext{Intervalo de tiempo en el video: 00:25:05--00:25:18.}
\begin{knowledgebox}{El valor de las trayectorias de búsqueda}
Las trayectorias de búsqueda permiten que el modelo «piense» antes de generar cada paso: puede reproducir la ruta de ejecución, cuantificar heuristics e incluso agregar restricciones estáticas, lo cual reduce la tasa de error durante la inference y hace viable la revisión visual.
\end{knowledgebox}

El expositor también enfatizó que, en los experimentos reales (Soang, Maze, entre otros), los datos de entrenamiento no son instrucciones vacías de un solo paso, sino el trace, los valores de costo y el plan generados por el solver. Al muestrear 64 planes sobre el trace y luego reordenarlos con heuristics, el modelo aprende cadenas de razonamiento más largas con pocas muestras.

\subsection{Search Former y DU Forer}

Sobre esta base, el equipo propuso Search Former e introdujo el trace drop en DU Forer. Al eliminar aleatoriamente parte de los tokens de búsqueda (por ejemplo, ciertas estimaciones de cost o intermediate node), el modelo aprende a recuperar el núcleo de la solución incluso cuando faltan datos, de manera similar al entrenamiento contra interferencias self-supervised.
\begin{importantbox}{La lógica de entrenamiento de DU Forer}
1) primero se usa el solver para muestrear el trace completo; 2) se hace drop de tokens específicos por nivel (por ejemplo, cost, create, exploration); 3) se hace que el Transformer reconstruya el plan a partir del trace residual; 4) se actualizan el Transformer y la drop layer mediante retropropagación del loss.
\end{importantbox}
\begin{enumerate}
    \item \textbf{Level 0:} Se conserva el trace completo como salida de línea base.
    \item \textbf{Level 1:} Se elimina el create cost y solo se conserva el path structure.
    \item \textbf{Level 2:} Se eliminan algunas acciones de exploración, lo que refuerza el modelado de los puntos de inflexión clave.
    \item \textbf{Level 3:} Solo se conserva el plan token final, y el modelo debe reconstruir por sí mismo el trace completo.
\end{enumerate}
\begin{figure}[H]
\centering
\includegraphics[width=\textwidth]{frame-dspy-architecture.jpg}
\caption{Arquitectura de DU Forer: Search Former genera el trace, que luego pasa por la drop layer y el policy head para producir el plan. \protect\footnotemark}
\end{figure}
\footnotetext{Intervalo de tiempo en el video: 00:32:25--00:32:45.}
\begin{warningbox}{Longitud del trace y equilibrio de costos}
Un trace largo aporta más contexto, pero también exige mayor longitud de tokens, memoria de GPU y sobrecarga de análisis; en la práctica es necesario controlar la estrategia de drop para evitar secuencias que sean matemáticamente completas pero inutilizables durante la inference.
\end{warningbox}

\subsection{Resumen de la sección}

Usar el trace de búsqueda como objetivo de aprendizaje, y regular la densidad de información mediante la estrategia de drop de DU Forer, permite que incluso los modelos pequeños capturen con eficiencia la calidad de la planificación.

\section{Diversidad de trace y estrategias de muestreo}

Para que trace mejorara el rendimiento sin dejar de ser interpretable, el equipo diseñó un pipeline de muestreo en varias etapas: primero el solver produce 64 candidate trace, y luego una puntuación heuristic (cost, length, constraint violation) decide cuáles trace vale más la pena conservar. Cada trace registra cost, exploration depth y plan segment; estos metadata también se usan para entrenar la supervisión de diversity.
\begin{enumerate}
    \item \textbf{Generación de trace:} se usa solver/Beam search para construir un trace global y registrar el intermediate cost.
    \item \textbf{Filtrado de trace:} se filtran los trace de baja calidad según violation, length y user goal.
    \item \textbf{Muestreo de trace:} se muestrean varios plan con el objetivo de la diversidad, para garantizar que se cubran distintas estrategias durante el entrenamiento y la inferencia.
    \item \textbf{Reordenamiento de trace:} la sequence que produce el Transformer se valida mediante un judge y luego se alinea con el original solver trace, para formar nuevos datos de entrenamiento.
\end{enumerate}
\begin{knowledgebox}{Los beneficios de la diversidad}
Esta estrategia de muestreo hace que el entrenamiento ya no dependa de un solo optimal plan, sino que el modelo aprenda la manera de pensar de «buscar varias soluciones factibles y elegir la mejor»; sobre todo en Soang, Maze y Nano optics, diversity mejora directamente el test-time performance.
\end{knowledgebox}

\subsection{Resumen de la sección}

Un muestreo diversificado de trace puede combinar la ventaja del deterministic solver con la flexibilidad del probabilistic planner, conservando la corrección del plan a la vez que aporta nuevos candidate, de manera que el modelo de search trace pueda seguir produciendo soluciones robustas en entornos unseen.

\section{Interacción de varios turnos y diseño de la constitución}

El segundo paso de Compound AI es la recolección de datos mediante diálogo de varios turnos. Cuando el usuario solo da una intención vaga, el sistema necesita preguntar los detalles, mantenerse confiable y dar instrucciones estructuradas dentro de un número limitado de turnos.

\subsection{APAC Constitution}

Ellos abstrajeron esta necesidad en la APAC Constitution: las cuatro dimensiones de Accuracy, Proactivity, Adaptivity y Credibility, cada una cuantificada mediante conductas concretas, por ejemplo si la pregunta se enfoca en la información clave real o si evita repetir preguntas irrelevantes varias veces.
\begin{figure}[H]
\centering
\includegraphics[width=\textwidth]{frame-apac-constitution.jpg}
\caption{APAC Constitution responde al usuario en varias dimensiones para lograr preguntas proactivas y alta confiabilidad. \protect\footnotemark}
\end{figure}
\footnotetext{Intervalo de tiempo en el video: 00:16:08--00:16:22.}
\begin{importantbox}{Lecciones de la constitución APAC}
1) Accuracy: los detalles deben caer en los campos clave que el usuario necesita; 2) Proactivity: iniciar preguntas de valor por cuenta propia; 3) Adaptivity: ajustar el estilo según el traveler persona; 4) Credibility: controlar la hallucination y evitar contradicciones evidentes. APAC hace posible optimizar estos objetivos al mismo tiempo.
\end{importantbox}
\textit{``Agent Constitution makes it possible to improve the agent performance along different axes simultaneously with very simple techniques.''}

Un ejemplo real de conversación es: cuando los travelers dicen ``I want to go to Hawaii'', el agent primero confirma el presupuesto, la preferencia (money vs experience) y si quiere varias ciudades, después pregunta la ventana de tiempo y los medios de transporte aceptables, y en cada turno actualiza la información recabada en el Json schema. Este pipeline de preguntas progresivas con validación inmediata está impulsado por los cuatro objetivos de APAC, lo que garantiza que las preguntas tengan valor sin desviarse de la intención inicial del usuario.

\subsection{Preguntas proactivas y evaluación}

Ellos simularon 50 traveler persona, cada uno con preferencias distintas y una hidden spreadsheet; el agent completa la entrada Json mediante varios turnos de diálogo y después se hace DPO fine-tune para optimizar la alineación con el ground truth. El experimento también usó al propio agent como judge, sin depender ya de una calificación humana durante todo el proceso, lo que aumenta la eficiencia de los datos.
\begin{knowledgebox}{Métricas de las preguntas de varios turnos}
Las evaluaciones clave incluyen: 1) overall recall (si se recabó toda la información); 2) critical recall (la corrección de los campos más importantes); 3) agentic score (la proactividad y la coherencia a lo largo de varios turnos de la respuesta); 4) DPO reward (la alignment con el judge).
\end{knowledgebox}
\begin{warningbox}{Evitar la trampa de un persona superficial}
Las preguntas proactivas deben subordinarse al objetivo: preguntar en exceso desperdicia turnos, y apegarse en exceso a una plantilla fija ignora fácilmente las particularidades del persona. El modelo debe encontrar el equilibrio entre accuracy y efficiency.
\end{warningbox}
\begin{knowledgebox}{Heurísticas para preguntar de forma proactiva}
1) Dar seguimiento a la información faltante: en cada turno ordenar las preguntas candidatas según la importancia del campo; 2) medir el cost-benefit: si una pregunta trae un valor alto pero un costo alto, se puede dividir en dos preguntas más pequeñas; 3) usar el persona para enfocar las diferencias: al tipo sensible al presupuesto preguntarle rápido por el costo, al orientado a la experiencia enfatizarle las actividades; 4) limitar los turnos: si se pasa de 4 turnos, considerar resumir lo que ya se sabe.
\end{knowledgebox}

\subsection{Resumen de la sección}

La constitución APAC le da a las preguntas de varios turnos control y la posibilidad de medirse, el agent judge junto con DPO forman un ciclo cerrado, lo que hace viable el fine-tuning autosupervisado como esquema de entrenamiento.

\section{Alineación humano-máquina y evaluación}

El expositor propone además la idea del ``agent judge'': en lugar de depender del etiquetado manual de todo el trace, se usa un agent optimizado por preferencias como judge, que califica el plan de cada interacción, y luego esas señales de preferencia se retroalimentan al proceso de generation.
\begin{table}[H]
\centering
\begin{tabular}{lp{6cm}}
\toprule
Métrica & Descripción \\
\midrule
Accuracy & Si los campos clave (destino, presupuesto, tiempo) son correctos y consistentes \\
Proactivity & Si formula suficientes preguntas dentro de un número limitado de turnos para completar el contexto \\
Credibility & Autoconsistencia y frecuencia de hallucination; el judge penaliza las respuestas que se contradicen a sí mismas \\
Diversity & Cantidad de planes en la conversación de varios turnos y cobertura de distintas estrategias \\
\bottomrule
\end{tabular}
\caption{Métricas de evaluación de agent judge}
\end{table}

\begin{knowledgebox}{Agent Judge+Human Alignment}
La salida del agent judge se somete a un soft alignment con proxies de human preference (como plan optimality y constraint compliance), y mediante DPO se retropropaga para que el generator converja más rápido hacia el comportamiento aprobado por el usuario.
\end{knowledgebox}

\subsection{Resumen de la sección}

El agent judge ofrece un bucle de retroalimentación de alta frecuencia: el reward de cada conversación se usa tanto para la actualización DPO del propio judge como para calibrar el planner, formando un ciclo cerrado de metrics + training.

\section{Slide y estrategia de evidencia visual}

\subsection{Archivo estructurado slide-first}
El expositor, al narrar el proceso de optimización de trace, agent y DSPy, se apoya siempre en imágenes de slide estructuradas para resaltar el pipeline completo. En el proceso de edición seguimos el principio slide-first: primero capturamos los slides estáticos (como las tres categorías de estrategias, la Constitución APAC, la arquitectura de DSPy y la latent optimization), y después, donde sea necesario, agregamos frames dinámicos, usando evidencia visual para anclar cada tramo de la lógica didáctica. La siguiente tabla enumera el material visual principal que se usa hasta ahora y su valor didáctico.
\begin{table}[H]
\centering
\begin{tabular}{lp{8cm}}
\toprule
Slide/Frame & Contenido principal y uso \\
\midrule
frame-solution-categories.jpg & El slide de «capas» donde aparecen las tres estrategias (Scaling/Hybrid/Compound AI), que ayuda al lector a entender el trade-off de costo/control entre los distintos métodos (00:09:10--00:09:18). \\
frame-apac-constitution.jpg & La tabla de las cuatro dimensiones de la APAC Constitution, que explica cómo se vinculan las métricas discriminativas con la estrategia de preguntas del agent (00:16:08--00:16:22). \\
frame-dspy-architecture.jpg & El diagrama de arquitectura de DSPy, que muestra el ciclo cerrado entre latent cost, solver y plan (00:32:25--00:32:45). \\
frame-latent-opt.jpg & El diagrama diferenciable que describe el pipeline de latent optimization, usado para explicar la iteración de dos niveles de Green Descent (00:45:05--00:45:25). \\
\bottomrule
\end{tabular}
\caption{Slides/frames usados en esta sección y su valor didáctico respectivo}
\end{table}
\begin{knowledgebox}{Sinergia entre slide y frame}
Los slides se usan para presentar fórmulas estructuradas, hojas de ruta y tablas de métricas; los frames se usan para capturar el estado del expositor mientras explica, el movimiento de un diagrama esquemático o un ejemplo numerical llamativo. Mientras haya un slide disponible, se prioriza el slide, y el frame sirve como complemento dinámico o para llenar puntos ciegos de la explicación; juntos sostienen una narrativa didáctica continua.
\end{knowledgebox}

\subsection{Notas de tiempo de los frames clave}
Cada ilustración registra estrictamente su intervalo de tiempo correspondiente, colocado en el pie de figura o en el footer, para que el lector pueda usar el subtitle y volver a los detalles. Por ejemplo, trace curve, la estructura APAC y el latent optimizer se marcan, cada uno en su propia subsección, con el formato \texttt{00:xx:xx--00:yy:yy}. Si en el futuro se agregan nuevos slides, es indispensable primero comparar con los subtitles para encontrar el frame más completo que se muestre, y anotar con claridad el intervalo de tiempo.
\begin{warningbox}{La marca de tiempo de la imagen no se puede omitir}
La falta de una marca de tiempo impide que el lector confirme el origen de la visual evidence, sobre todo cuando varias imágenes contiguas tratan temas distintos. Cada figure debe dar el intervalo de subtítulos en la misma página, y de ser necesario, repetirlo en el caption.
\end{warningbox}

\subsection{Resumen de la sección}
La forma de organización slide-first le da al pipeline abstracto anclas visuales, y junto con los ejemplos dinámicos del frame y las marcas de tiempo precisas, permite que el lector, al leer el texto, regrese rápidamente a la escena original de la explicación en el video.

\section{Optimización compleja y codificación de restricciones en DSPy}

La última parte describe cómo el marco DSPy maneja las restricciones no lineales difíciles de expresar en la optimización combinatoria, como el diseño birefringent en nanophotonics.

\subsection{Linealización latente \& Green Descent}

Para las restricciones complejas, el equipo primero predice un cost vector \(C\), luego convierte el problema original en un subproblema lineal y, finalmente, usa un solver para generar \(X^\star\). Todo el proceso forma un grafo de cómputo que va de la description \(Y\) a \(C(Y)\) y luego a la salida del solver; el loss es la diferencia respecto al objetivo original en el dominio temporal, y se retropropaga mediante gradient descent para entrenar la representación de \(C\).
\begin{figure}[H]
\centering
\includegraphics[width=\textwidth]{frame-latent-opt.jpg}
\caption{Grafo de cómputo en DSPy que deriva \(C\) a partir de la descripción \(Y\) y luego llama al solver para producir \(X^\star\). \protect\footnotemark}
\end{figure}
\footnotetext{Intervalo de tiempo del video: 00:45:05--00:45:25.}
\begin{knowledgebox}{Latent Optimization Pipeline}
1) A partir de la descripción se predice el latent cost \(C\); 2) con \(C\) fijo, se resuelve de manera aproximada con un solver diferenciable; 3) la salida se reinserta en el loss físico; 4) se entrena \(C\) con descenso de gradiente. Esta resolución inversa, veloz como un relámpago, es más estable que resolver directamente el problema no lineal original.
\end{knowledgebox}
\begin{importantbox}{Pasos jerárquicos de Green Descent}
Primero se usa un solver ligero para resolver analíticamente el subproblema linealizado y generar \(X^\star\); después la salida pasa por backprop de regreso al latent cost, y con descenso de gradiente se ajusta \(C\). Esta iteración de dos niveles (solver + actualización de gradiente) evita numéricamente las oscilaciones que produce ascender directamente por el espacio no lineal original.
\end{importantbox}

\subsection{Restricciones ópticas y de fabricación}

En el diseño óptico de una cuadrícula de 80\u00d780, cada punto de la cuadrícula solo puede tomar el valor 0/1, y el proceso de fabricación exige que no aparezcan puntos aislados y que se respete el tamaño del pincel. También hay que considerar restricciones multidimensionales como la respuesta en frecuencia y las trayectorias de interferencia, por lo que el problema se descompone y se valida en varios benchmark (beam splitter, wavelength multiplexer).
\begin{warningbox}{Trampas ocultas de las restricciones de fabricación}
Las restricciones no solo tienen forma algebraica, sino que también incluyen la fabricabilidad de ingeniería: las restricciones lineales derivadas del tamaño del pincel, los píxeles aislados inviables y las limitaciones topológicas de las trayectorias de la luz. Cualquier optimización que ignore esto resultará inservible en la etapa de fabricación.
\end{warningbox}

\subsection{Resumen de la sección}

Mediante la combinación de latent linearization y la retroescritura del solver, DSPy logra que el aprendizaje de restricciones complejas pase de un ``espacio de búsqueda infinito'' a ``parámetros diferenciables'', y a la vez incorpora las restricciones de fabricación al sistema de evaluación.

\section{Ingeniería y validación}

Para validar Compound AI en la ingeniería real, el expositor presentó varios benchmarks: el juego Soang, la búsqueda de rutas en Maze y el diseño de rejillas de 80\u00d780 para óptica a escala nanométrica. En cada categoría se compara un agent (o un latent solver) entrenado con search trace y DPO contra soluciones puras y modelos generativos puros.

\subsection{Puntos de referencia de Soang y Maze}

Soang es un box-pushing game en el que solo se puede empujar hacia adelante: hay que planear bien el orden de las jugadas, porque en cuanto se empuja un box a una esquina ya no se puede regresar. En los experimentos, el Solution-only baseline necesitó cerca de 175 millones de parámetros y muestras de entrenamiento del orden del millón para tener éxito solo de manera ocasional; en cambio, Search Former+DU Forer alcanzó 80\% de exactitud con 15M de parámetros y 100,000 muestras.
El benchmark de Maze usa labyrinth de distintos tamaños: (1) mazes pequeños de sub 20 steps; (2) mapas grandes que se extienden hasta 40+ steps; (3) con obstáculos dinámicos incorporados. El modelo de search trace mostró una generalization más fuerte que el LLM puro, y en los Maze con muchas bifurcaciones, el search trace también ofrece un candidate pruning eficaz.
\begin{knowledgebox}{Resumen de resultados de ingeniería}
En la tarea de Soang, el trace dropout mejoró la robustness; en la tarea de Maze, el consistency check (un judge de autoverificación del agent) redujo de forma significativa la hallucination. Ambos benchmarks muestran que el modelo de trace, con 5 a 10 veces menos datos, logra igualar o incluso superar al solution-only baseline convolucional.
\end{knowledgebox}

\subsection{Benchmark de óptica a escala nanométrica}

En la tarea de diseño óptico, el objetivo es usar un binary pattern de una rejilla de 80\u00d780 para controlar la refracción de cada píxel. El pipeline de entrenamiento requiere: 1) simular la respuesta en frecuencia; 2) calcular el loss del beam splitter y del wavelength multiplexing; 3) tener en cuenta las restricciones del brush durante la fabricación.
El expositor enfatiza: solo incorporando estos indicadores al loss, el latent cost \(C\) puede aprender diseños que sean a la vez válidos y de alto desempeño. Además, el proceso de retroalimentar el solver hacia \(C\) permite que el modelo aprenda de manera automática a partir del failure trace, y así se transfiera entre distintos benchmarks.

\subsection{Resumen de la sección}

Estos benchmarks abarcan tanto el search al estilo del ajedrez como el diseño continuo de óptica, y demuestran que Compound AI puede, en un solo modelo, encadenar múltiples módulos funcionales en distintos dominios y completar el entrenamiento mediante señales de supervisión sencillas.

\section{Análisis de casos}

\subsection{Soang: de la búsqueda a DPO}

Soang es un juego de box-pushing en el que solo se puede empujar hacia adelante, y pone a prueba en gran medida la capacidad del sistema para planificar los estados futuros. El enfoque de Compound AI se divide en dos etapas: primero se usa el trace de búsqueda para aprender el orden estricto de las acciones, y después se usa DPO para que el agent judge dé retroalimentación y reduzca el action spirit. El flujo de entrenamiento es el siguiente:
\begin{enumerate}
    \item Se usa el solver para generar el trace completo, registrando cost, plan e intermediate state;
    \item Se etiqueta el trace (si cumple las constraints, si se completa en 4 rondas);
    \item Search Former aprende el trace; el agent judge muestrea varios plan y evalúa la fidelity;
    \item El fine-tuning con DPO hace que el agent judge se incline más hacia el plan óptimo.
\end{enumerate}
El resultado final es que Search Former + DU Forer, con 15M de parámetros y 100K trace, alcanza una tasa de éxito del 80\%, mientras que el baseline solution-only necesita 175M de parámetros y 1M de trace para llegar solo al 20\%.
\begin{table}[H]
\centering
\begin{tabular}{lrrl}
\toprule
Modelo & Parámetros & Muestras & Tasa de éxito \\
\midrule
Solution-only & 175M & 1M & 20\% \\
Search Former + DU Forer & 15M & 100K & 80\% \\
\bottomrule
\end{tabular}
\caption{Comparación de modelos en el benchmark de Soang}
\end{table}

\subsection{Nano optics: del latent cost a la fabricación}

En el diseño óptico de una malla de 80\u00d780, el trace de búsqueda se convierte en un trace de simulación: el agent observa la trayectoria de la luz, la respuesta en frecuencia y el interference pattern, y después comprime esta información en el latent cost \(C\). Todo el pipeline incluye: 1) generar la simulación del beam splitter y el wavelength multiplexer; 2) elegir el binary pattern que satisface la restricción de brush; 3) usar el solver para producir \(X^\star\); 4) propagar el loss hacia atrás hasta \(C\).
\begin{importantbox}{El ciclo cerrado del trace a la fabricación}
1) El trace conserva la información de la interferencia de la luz y de la brocha de fabricación; 2) el latent \(C\) extrae la descripción del entorno y a la vez controla al solver; 3) la retroalimentación del solver se usa para calibrar \(C\) y evitar generar puntos no fabricables; 4) el multi-task loss permite que el diseño se transfiera entre distintos benchmark.
\end{importantbox}
\begin{knowledgebox}{Enseñanzas de los casos}
Estos dos casos muestran en conjunto que Compound AI puede encadenar el trace de planificación, las preguntas y respuestas interactivas y el Latent solver, de modo que cada salida es a la vez la entrada de la siguiente etapa y un resultado intermedio auditable. Este modular pipeline facilita localizar las fallas y hacer iteraciones en pasos pequeños.
\end{knowledgebox}

\subsection{Resumen de la sección}

El análisis de casos destaca el flujo de datos en la práctica de la ingeniería: el search trace, el agent judge y el latent cost se conectan entre sí, y se validan repetidamente en distintas tareas para garantizar que la teoría abstracta pueda llevarse a la práctica.

\section{Recursos y transparencia de costos}

Yuandong Tian recuerda en varias ocasiones: el núcleo de un marco unificado no consiste en seguir apilando modelos grandes en un solo punto, sino en repartir recursos limitados entre la búsqueda, el muestreo de diálogos, la evaluación y la optimización, que son etapas acopladas entre sí.

\subsection{Comparación de costos de entrenamiento e inferencia}

\begin{table}[H]
\centering
\begin{tabular}{lrrr}
\toprule
Componente & Consumo principal & Cantidad típica de parámetros & Métrica observada \\
\midrule
Search Former & GPU hours × trace length & 10-20M & Trace fidelity, diversity coverage \\
APAC agent & Token cost × 4-6 rounds & 500M-1B & Dialogue throughput, judge reward \\
DSPy solver & Solver iterations × latents & 5-10M & Constraint slip rate, latents convergence \\
\bottomrule
\end{tabular}
\caption{Distribución de recursos de cada módulo de Compound AI}
\end{table}

\begin{importantbox}{Estrategia de asignación}
1) Ahorrar tokens en el search trace, pero conservar la información clave de costo; 2) en APAC, usar lightweight fine-tuning para iterar rápido el judgment; 3) en DSPy, dar seguimiento a los solver steps para evitar el gradient explosion. Esta combinación de recursos resulta más económica que un simple scaling.
\end{importantbox}

\subsection{Control de la latencia de inferencia}

En la inferencia, las tres etapas — trace generation → APAC dialogue → DSPy solver — se encadenan, y la latencia total queda limitada por la etapa más lenta. En el despliegue real se adoptan las siguientes estrategias:
\begin{itemize}
    \item Search trace precompute + cache heatmap;
    \item En APAC se fija un mínimo de 3 rondas; si se superan las 5 rondas se activa el fallback;
    \item El DSPy solver usa warm start, es decir, reutiliza la inicialización de \(C\) de la vez anterior.
\end{itemize}

\begin{knowledgebox}{Latency guardrails}
En el pipeline de inferencia se configuran \texttt{max\_trace\_steps}, \texttt{max\_dialog\_rounds} y \texttt{max\_solver\_iters}, para evitar que una etapa anómala arrastre a toda la cadena. El monitoreo debe incluir la longitud de la cola y la duración del rollout.
\end{knowledgebox}

\subsection{Resumen de la sección}

Descomponer los recursos (trace / agent / solver) permite medir con mayor precisión el costo y el beneficio de Compound AI, y al mismo tiempo establecer guardrails de latencia en la inferencia.
\section{Recomendaciones de despliegue y depuración del sistema}

Para desplegar este Compound AI completo en un producto real, es necesario establecer interfaces claras entre search, APAC y DSPy, y agregar monitoreo a cada resultado intermedio:
\begin{itemize}
    \item El Search trace output debe registrar los cost/constraint flags, para facilitar el debug;
    \item El multi-step question log de APAC debe conservar el intent del usuario más la agent ask list;
    \item El latent cost \(C\) y el solver output de DSPy deben reportarse juntos, para poder alinear distintos benchmark.
\end{itemize}

\begin{importantbox}{Puntos clave de depuración y despliegue}
1) Agregar monitor: en inference, establecer un sanity check para cada trace, question y plan; 2) agregar caché: el search trace suele repetirse, por lo que puede almacenarse en caché y reutilizarse; 3) agregar fallback: cuando el agent judge señale una violación de restricción, retroceder a la capa del hybrid solver para replanificar; 4) agregar control de versiones: cada entrenamiento DPO debe ir acompañado de la versión del modelo judge.
\end{importantbox}

\begin{knowledgebox}{Cómo mantener transparente la cadena}
Escribir un structured log (trace id, agent response, solver plan) en cada etapa del pipeline, y proporcionar para cada log un mapeo de \texttt{trace id -> user request}; de este modo, incluso si el sistema falla en producción, es posible reproducir con rapidez la causa del error.
\end{knowledgebox}

\subsection{Resumen de la sección}

Al desplegar Compound AI, es necesario que la salida de cada módulo sea observable y reproducible, y establecer un enlace de versiones entre judge, solver y trace, para poder ubicar con rapidez las fallas de razonamiento y ejecutar reversiones.

\section{Observabilidad y métricas de evaluación}

Para que Compound AI funcione de manera estable en la práctica de ingeniería, es necesario mapear el estado de ejecución de trace, agent y solver a un conjunto claro de métricas.

\subsection{Panel de métricas clave}

\begin{itemize}
    \item \textbf{Trace fidelity}: la similitud entre el trace y el plan que produce el solver (por ejemplo, el Jaccard overlap).
    \item \textbf{Judge reward curve}: la variación, ronda a ronda, de accuracy/proactivity/credibility del agent judge.
    \item \textbf{Constraint slip rate}: la frecuencia con la que aparece un constraint violation en cada plan.
    \item \textbf{Dialogue throughput}: el número de rondas que requiere el diálogo de preguntas y respuestas de APAC para completar un Json.
\end{itemize}

\begin{importantbox}{La ruta de datos detrás de las métricas}
Cada métrica debe vincularse a una fuente de datos rastreable: trace fidelity se obtiene de la search trace cache, judge reward curve se muestrea directamente del log stream del DPO judge, y además, del lado de inference, también hay que registrar el \texttt{trace\_id} del plan. Constraint slip rate requiere que las violations que devuelve el solver se alimenten a la observability db; Dialogue throughput, por su parte, debe sumar latency instrumentation (el judge posterior adjunta el número de rondas y el token count). Solo así el dashboard puede reflejar con precisión, en cuestión de segundos, el estado del sistema.
\end{importantbox}

\begin{knowledgebox}{Los umbrales de alerta de las métricas}
Se configuran guardrails: si el judge reward desciende durante tres rondas consecutivas, se activa un new DPO fine-tune; si trace fidelity < 0.7, se retrocede al deterministic solver; si constraint slip rate aumenta, se agrega un \texttt{penalty\_update}. Estos umbrales de alerta permiten que el equipo de operaciones detecte rápidamente el drift del modelo.
\end{knowledgebox}

\subsection{Confiabilidad a nivel de trace}

La confiabilidad del trace puede evaluarse mediante la cobertura: cada inference registra un trace ID y el backend calcula un coverage heatmap a partir de los trace metadata (e.g., depth, constraint violations), lo que garantiza que el conjunto de entrenamiento y el conjunto de prueba tengan suficientes muestras en cada region.

\begin{warningbox}{No ignores el trace skew}
Los Long tail trace (rutas muy poco frecuentes, ya sea extremadamente largas o con large cost) se ignoran fácilmente durante el entrenamiento, pero suelen ser también una fuente de failure. Debe construirse un trace-skew dashboard para volver a muestrear estos casos a tiempo.
\end{warningbox}

\subsection{Resumen de la sección}

Contar con métricas a nivel de dashboard permite hacer observable el pipeline abstracto de Compound AI: en cuanto algún módulo produce una salida anómala, es posible ubicarla rápidamente siguiendo la ruta trace → judge → solver, y corregir el rumbo de manera automática cuando se activa un guardrail.
\section{Visualización y monitoreo de ingeniería}

\subsection{Diseño del dashboard}
Cada módulo de Compound AI debe tener un tablero visible en el panel operativo: la trace fidelity curve permite comparaciones horizontales, la judge reward curve y la user satisfaction signal se combinan en una vista de dos en uno, y la salida del solver destaca la constraint violation rate en otro panel. La siguiente tabla enumera los tableros clave y la frecuencia de actualización recomendada.
\begin{table}[H]
\centering
\begin{tabular}{lp{8cm}r}
\toprule
Tablero & Métricas incluidas & Frecuencia de actualización recomendada \\
\midrule
Trace health & fidelity, cost variance, skew coverage & 1 minuto \\
Judge stability & distribución de accuracy/proactivity/credibility, NPS proxy scores & 30 segundos \\
Constraint guard & solver violation rate, cumplimiento del proceso de brush & 10 segundos \\
\bottomrule
\end{tabular}
\caption{Principales dashboards de Compound AI y estrategia de actualización}
\end{table}
\begin{knowledgebox}{Cadencia de retroalimentación de la visualización}
Estratifica la frecuencia de actualización del monitoreo de los distintos módulos según su criticality, para evitar que actualizar todos los indicadores al mismo tiempo sature el monitoreo; a la vez, usa \texttt{trace\_id} como dimensión común para enlazar en el dashboard la información de las tres secciones trace, judge y solver, lo que ayuda a los ingenieros a investigar entre módulos.
\end{knowledgebox}

\subsection{Reproducción de Trace y Plan}
La transparencia real proviene de la capacidad de reproducción de trace/plan: después de ejecutar cada plan se escribe en \texttt{trace\_id.json}, que incluye la puntuación heuristic, las violaciones de restricciones y el judge reward, lo que permite reproducir paso a paso, en la página web de operaciones, el search trace, el judge dialogue y el solver output. Para que cada reproducción tenga una impresión visual, registramos la marca de tiempo de los frame clave en el mismo log, de manera que al reproducir se pueda cargar la captura de la slide o la imagen del frame correspondiente.
\begin{knowledgebox}{La reproducción mejora la velocidad de auditoría}
La reproducción de un trace no solo permite reconstruir el decision path, sino que también deja que el engineer vea directamente la visual evidence alineada con la slide. Por ejemplo, la slide de APAC del segmento 00:16:08--00:16:22 se usa en la reproducción como anotación del judge trigger; el diagrama de arquitectura de DSPy del segmento 00:32:25--00:32:45 se usa para explicar el cálculo step-by-step del solver.
\end{knowledgebox}
\begin{warningbox}{Los datos de reproducción son indispensables}
La falta de alineación temporal entre trace y frame retrasa el progreso de la investigación. Es indispensable escribir en el registro, al mismo tiempo, \texttt{trace\_id}, \texttt{frame\_path} y \texttt{timestamp}; de lo contrario, al reproducir solo se podrá contar con una comparación imprecisa y el juicio operativo se distorsionará.
\end{warningbox}

\subsection{Resumen de la sección}
El punto central del monitoreo de ingeniería es combinar metrics, visual evidence y la lógica de reproducción en una cadena triple consultable, de manera que, cuando ocurra un evento, se pueda localizar rápidamente el trace frame, el judge dialog y el solver plan correspondientes, y así responder con precisión a las fallas.

\section{Repaso y reflexión de la sección}

\paragraph{}
Los primeros 15 minutos de la Lecture 05 delimitaron el dominio de planning domain mediante curvas de costo computacional y desempeño: la ``smooth curve'' del LLM ya mostraba rendimientos marginales decrecientes del scaling, por lo que era necesario buscar una ruta de solución más estructurada.

\paragraph{}
Search Former y DU Forer convierten el search trace en una secuencia con supervision, y usan dropout para simular la información de búsqueda incompleta que existe en la realidad. Los benchmarks de Soang y Maze mostraron que el modelo de trace puede, con un modelo pequeño, superar al baseline solution-only en generalization, y ofrecer recomendaciones más consistentes para el multi-turn agent.

\paragraph{}
La parte de interacción de múltiples turnos forma, mediante la APAC Constitution y el DPO loop del agent judge, un ciclo autosupervisado que enfatiza la medición en cuatro dimensiones y una estrategia de preguntas diferenciadas por persona, lo que permite que el agent produzca un plan estructurado dentro de un número limitado de turnos y reduzca de forma notable la frecuencia de alineación manual.

\paragraph{}
El marco DSPy utiliza el latent cost \(C\) como puente para enlazar el search trace, la salida de APAC y el solver, reduciendo la dimensionalidad del problema de combinatorial constraint hasta obtener una trayectoria de optimización diferenciable. El Nano optics benchmark demostró que este pipeline puede atender al mismo tiempo el desempeño y la validez del proceso de fabricación.

\paragraph{}
Lo más destacado de esta clase es que enlaza search, agent y solver en un solo flujo de información, en lugar de tratarlos como módulos independientes: cada salida incluye un time-stamped trace, un judge reward y un solver plan, y en conjunto forman un audit log completo.

\subsection{Resumen de la sección}
Al repasar esta sección se observa que, desde el search trace hasta el multi-turn agent y de ahí al DSPy latent solver, toda la arquitectura gira en torno a la interpretabilidad, la supervisión estructurada y la observabilidad de la ingeniería, formando un bucle cerrado reproducible y depurable.

\section{Síntesis y ampliación}

Esta clase mostró en una hora, mediante la ruta de tres pasos de Compound AI: 1) capturar y comprimir las trayectorias de búsqueda, 2) usar la Constitución APAC para gobernar las preguntas y respuestas de múltiples turnos, y 3) aprender restricciones complejas mediante DSPy, el ciclo completo desde el modelado hasta la implementación en ingeniería.

Estas tres partes no son independientes entre sí: el modelo de trace se usa para guiar el flujo de preguntas y respuestas del multi-turn agent, el Json estructurado que genera APAC es a su vez la entrada de restricciones del latent solver en DSPy, y el plan de acción final se envía de vuelta al search trace para una verificación de consistencia, formando un anillo de control auditable.

\begin{table}[H]
\centering
\begin{tabular}{lp{5cm}p{6cm}}
\toprule
Tema & Mecanismo clave & Valor didáctico \\
\midrule
Trayectoria de búsqueda & Trace-aware Transformer + DU Forer drop & Mejora la data/parameter efficiency y refuerza la interpretabilidad \\
Interacción de múltiples turnos & APAC Constitution + agent judge/DPO & Las preguntas de múltiples turnos poseen accuracy/proactivity/credibility \\
Restricciones complejas & Latent cost \(C(Y)\) + solver + green descent & Incorpora retroalimentación diferenciable en escenarios no lineales, sin dejar de lado las restricciones de manufactura \\
\bottomrule
\end{tabular}
\caption{Componentes clave y valor del método de tres pasos de Compound AI}
\end{table}

\subsection{Preguntas de investigación y extensiones}

\paragraph{Gestión de versiones de la capa de entrenador de IA}
Cuando el search trace, el agent APAC y el solver de DSPy iteran al mismo tiempo, es fácil que surja una «desincronización entre los parámetros del trace y el juicio del agent». Se recomienda empaquetar juntos el ID del trace, la versión del modelo de agent judge y la solver configuration en un registro de release, y ejecutar una prueba automatizada antes de cada despliegue, para asegurar que la nueva versión del judge todavía pueda reconocer los trace antiguos y evitar una unstable inference.

\paragraph{Generalización entre dominios y calibración del judge}
El agent judge actual destaca en tareas de travel intent, pero en ámbitos de alto riesgo como la medicina o las finanzas su comprensión del goal podría ser completamente distinta. Una estrategia de extensión viable es entrenar un mixture-of-experts para el judge o introducir un domain embedding que le permita adaptarse a un new intent, y a la vez alinear el judge reward con la human satisfaction signal (NPS, tasa de clientes recurrentes), en lugar de depender únicamente de la plan optimality.

\paragraph{Fusión multimodal y materiales de apoyo}
El expositor mencionó que «los slides, notebooks y photos» son herramientas poderosas para reforzar la explicación. Idealmente, el search trace debería poder referenciar las tablas/fórmulas de los slides (como el sequencing graph de Soang); cuando aparezca un diagram clave en el video, capturar automáticamente el fotograma e insertarlo en la note, manteniendo una doble evidencia visual y textual.

\begin{itemize}
    \item ¿Cómo comprimir aún más la cantidad de tokens del trace para reducir el costo de inferencia, manteniendo al mismo tiempo la interpretabilidad de múltiples turnos?
    \item ¿Cómo puede transferirse la capacidad de generalización del agent judge a ámbitos de alto riesgo como la medicina o las finanzas, y no solo al travel intent?
    \item ¿Pueden las restricciones de manufactura diferenciadas abstraerse automáticamente en una latent representation, de modo que se transfieran sin fricciones a nuevas industrias (por ejemplo, el diseño de chips)?
\end{itemize}

Además de las preguntas teóricas, este Compound AI también debe consolidarse de manera continua en production: mantener el audit log de trace/plan, alinear el DPO reward con la satisfacción del usuario, e iterar el modelo de judge en distintas region, para así asegurar la robustez del sistema a largo plazo.

\subsection{Lecturas adicionales}
\begin{itemize}
    \item «Beyond A*: Better Planning with Transformers via Search Dynamics Bootstrapping», útil para entender cómo las trayectorias de búsqueda mejoran el desempeño de la planificación.
    \item «Dualformer: Controllable Fast and Slow Thinking by Learning with Randomized Reasoning Traces», que corresponde al trace drop y a la velocidad de razonamiento controlable.
    \item «Composing Global Optimizers to Reasoning Tasks via Algebraic Objects in Neural Nets», que analiza la optimización global composicional en redes neuronales.
    \item «SurCo: Learning Linear Surrogates For Combinatorial Nonlinear Optimization Problems», que muestra una ruta de aprendizaje de surrogate para problemas de optimización combinatoria no lineal.
    \item APAC Constitution (2024), libro blanco interno que define en detalle los métodos de evaluación de Accuracy/Proactivity/Adaptivity/Credibility.
\end{itemize}

\subsection{Resumen de la sección}
El método de tres pasos de Compound AI no es un experimento aislado, sino una combinación orientada a production: fusiona los visual slides, los time-stamped frames y los trace logs en un mismo pipeline, para que la audiencia no solo entienda los principios del método, sino que también pueda dar seguimiento a la implementación concreta en ingeniería.

\end{document}
